%% sections/k-currency-matrix.tex - H. R. 9510 Bill v4.0 (full bill) - Appendix K
%% Genre: audit / currency matrix (verification matrix + sign-off block). Source
%% deliverable 11. Table 10. Gray-scale Mermaid figure (TikZ): Figure 9 (prompt
%% evolution beside bill-version evolution). Single hyphens only; section symbol
%% only in tables.
\billsec{APPENDIX K.}{CURRENCY AND CROSS-REFERENCE MATRIX (NON-OPERATIVE).}\phantomsection\label{toc:aK}

\uscprov{1}{\textit{Independent research draft. Not an enacted law, not pending
legislation, and not legal advice; not endorsed by the FDA, HHS, the Office of the Law
Revision Counsel, or any Member of Congress. This appendix is a verification record,
not an assertion that the law cannot change.}}

\uscprov{1}{\textbf{Scope.} This matrix records the currency and cross-reference audit
for H. R. 9510 (illustrative number; the Clerk assigns the real number at
introduction), 119th Congress, 2d Session, the Verification Before Generation in
Physical AI Oncology Trials Act of 2026. The bill amends the Federal Food, Drug, and
Cosmetic Act (FD\&C Act), 21 U.S.C. 301 et seq., as reproduced from Title 21, United
States Code, current through Public Law 119-93 as of the draft date, June 7, 2026. The
operative mechanism is a new section 515D of that Act, codified at 21 U.S.C. 360e-5,
inserted immediately after section 515C (21 U.S.C. 360e-4)
\cite{pl-119-93, olrc-usc21, usc-360e4}.}

\uscprov{1}{\textbf{Purpose and method.} The draft was frozen and currency-checked
through Public Law 119-93, and the legal crosswalk that anchors its citations is
current through May 31, 2026. This matrix performs the currency task that the
next-steps memo identifies as a precondition to introduction: re-verify, immediately
before the bill is introduced, that no newer public law has changed Title 21 since
Public Law 119-93; that the target section number 515D, codified at 21 U.S.C. 360e-5,
is still unused; and that every cross-reference and each of the ten conforming
amendments still lands on existing, correctly numbered law. The method is documentary
and reproducible. Each row names the authority it depends on, the status confirmed as
of June 7, 2026, and the precise action to repeat at introduction. The controlling
currency sources are the Office of the Law Revision Counsel (OLRC) United States Code,
the text of Public Law 119-93, and the legal crosswalk current through May 31, 2026
\cite{olrc-usc21, vvuq05-next-steps}.}

\uscprov{1}{Title 21 is amended frequently, so the value of the matrix is that it tells
a sponsoring office exactly which anchors to re-pull, and what a clean result looks
like, on the day the measure is filed. Table 10 is the verification matrix: each row
gives the citation, what to verify, and its status as of the draft date.}

{\footnotesize
\begin{xltabular}{\textwidth}{@{}L{2.6cm} Y L{2.4cm}@{}}
\hline\noalign{\vskip 0.04cm}
\textbf{Citation} & \textbf{What to verify} & \textbf{Status (June 7, 2026)}\\
\hline\noalign{\vskip 0.04cm}
\endfirsthead
\hline\noalign{\vskip 0.04cm}
\textbf{Citation} & \textbf{What to verify} & \textbf{Status (June 7, 2026)}\\
\hline\noalign{\vskip 0.04cm}
\endhead
\hline\noalign{\vskip 0.04cm}
\endlastfoot
Title 21, chapter 9 & No public law after Public Law 119-93 has amended the FD\&C Act
chapter (21 U.S.C. 301 et seq.) & Current through Public Law 119-93; no later
amendment reflected\\
\hline\noalign{\vskip 0.04cm}
\S~301 (note) & The new Short Title of 2026 Amendment note is additive and does not
collide with the existing 2025 note & 2025 note present; 2026 note open\\
\hline\noalign{\vskip 0.04cm}
\S~331 & Last lettered prohibited act is (iii), so the new act lands as (jjj) & Last
act confirmed at (iii); (jjj) open\\
\hline\noalign{\vskip 0.04cm}
\S~351 & Subsections run through (j), so the new act lands as (k) & (h), (i), (j)
present; (k) open\\
\hline\noalign{\vskip 0.04cm}
\S~355g & Subsections (a) and (b), so the new act lands as (g) & (a), (b) confirmed;
(g) open\\
\hline\noalign{\vskip 0.04cm}
\S~360e(c)(1) & Clauses run through (H), so the new clause lands as (I) & (G), (H)
confirmed; (I) open\\
\hline\noalign{\vskip 0.04cm}
\S~360c & Subsections run through (k), so the new act lands as (l) & Through (k)
confirmed; (l) open\\
\hline\noalign{\vskip 0.04cm}
\S~360e-4 & Subsections run through (c), so the new act lands as (d) & Through (c)
confirmed; (d) open\\
\hline\noalign{\vskip 0.04cm}
\S~360k & Subsections (a) and (b), so the savings clause lands as (c) & (a), (b)
confirmed; (c) open\\
\hline\noalign{\vskip 0.04cm}
\S~360e-5 (\S~515D) & The target number is unassigned in both the Code and the FD\&C
Act table of contents & Confirmed unused; well placed after section 515C\\
\hline\noalign{\vskip 0.04cm}
Legal crosswalk & The crosswalk anchoring the citations has not been revised in a way
that moves an anchor & Current through May 31, 2026\\
\hline
\end{xltabular}}
\vspace{-0.8cm}
\figcaption{Table 10. The currency and cross-reference verification matrix: each
citation, what to verify, and its status as of the draft date. A clean result at
introduction is every row confirmed and section 360e-5 still unused.}

\uscprov{1}{\textbf{Target-number availability.} The Act needs one vacant section
number coordinated across two systems: the FD\&C Act's own section numbering and the
United States Code classification. The chosen pair is FD\&C Act section 515D,
classified to 21 U.S.C. 360e-5. The premarket-approval cluster runs section 515 (21
U.S.C. 360e), section 515A (21 U.S.C. 360e-1), section 515B (21 U.S.C. 360e-2), and
section 515C (21 U.S.C. 360e-4), with the predetermined change control plan authority
at section 515C. Inserting section 515D, codified at 21 U.S.C. 360e-5, immediately
after section 515C places the new verification authority next to the change-control
authority it most closely resembles, and the suffix is the next open slot in the
section 360e family. As of June 7, 2026, neither section 515D nor 21 U.S.C. 360e-5 is
assigned in Title 21. To re-verify at introduction, search the OLRC United States Code
for 360e-5 and the FD\&C Act table of contents for 515D; both must return no existing
section. If either has been taken by an intervening law, renumber to the next open
suffix and re-run the conforming-amendment matrix, because every conforming amendment
cross-references section 515D \cite{usc-360e4, olrc-usc21}.}

\uscprov{1}{\textbf{Section-number disambiguation.} A recurring drafting hazard in
Title 21 device law is the near-identical pair section 360(k) and section 360k. They
are different provisions and must not be conflated. Section 360(k) is subsection (k) of
section 360 (FD\&C Act section 510(k)), the premarket-notification authority, on which
the section 510(k) conforming amendment and the change-control cross-references in
section 515C rely. Section 360k is a separate, freestanding section (FD\&C Act section
521) on State and local requirements respecting devices, that is, device preemption,
to which the Act adds a new savings clause. Every operative use of 510(k) in the bill
cites 21 U.S.C. 360(k) with the parenthetical, and every use of the preemption section
cites 21 U.S.C. 360k without parentheses \cite{usc-360, usc-360k}.}

\uscprov{1}{\textbf{Build lineage.} The currency record sits within a larger build
history. Figure 9 places the prompt evolution beside the bill-version evolution: a
sequence of prompts (next-steps, figures-bill, draft-bill, full-bill, final-bill) runs
alongside the bill versions (v1.0 from VVUQ-03, v2.0 from VVUQ-04, v3.0 from VVUQ-05,
and v4.0, this build), and the final-bill prompt produces this v4.0 build, the fourth
bill perspective \cite{vvuq05-next-steps, kawchak2026vvuq05bill}.}

\begin{mermaidfig}
\node[mmstep,text width=24mm] (p1) at (0,0) {next-steps};
\node[mmstep,text width=24mm] (p2) at (0,-1.2) {figures-bill};
\node[mmstep,text width=24mm] (p3) at (0,-2.4) {draft-bill};
\node[mmstep,text width=24mm] (p4) at (0,-3.6) {full-bill};
\node[mmstep,text width=24mm] (p5) at (0,-4.8) {final-bill};
\draw[mmedge](p1)--(p2);\draw[mmedge](p2)--(p3);\draw[mmedge](p3)--(p4);\draw[mmedge](p4)--(p5);
\node[mmin,text width=24mm] (b1) at (4.4,0) {v1.0 VVUQ-03};
\node[mmstep,text width=24mm] (b2) at (4.4,-1.2) {v2.0 VVUQ-04};
\node[mmdec,text width=24mm] (b3) at (4.4,-2.4) {v3.0 VVUQ-05};
\node[mmgoal,text width=24mm] (b4) at (4.4,-3.6) {v4.0 this build};
\draw[mmedge](b1)--(b2);\draw[mmedge](b2)--(b3);\draw[mmedge](b3)--(b4);
\draw[mmedged] (p5) -- (b4);
\end{mermaidfig}
\vspace{-0.3cm}
\figcaption{Figure 9. Prompt evolution beside bill-version evolution: the final-bill
prompt produces this v4.0 build, the fourth bill perspective.}

\uscprov{1}{\textbf{Re-verification checklist before introduction.} Each item is a
clean result the verifier of record confirms on the day the measure is filed.}

\begin{enumerate}[leftmargin=2.0em,itemsep=2pt,topsep=2pt,label=(\arabic*)]
\item Pull the OLRC United States Code currency banner for Title 21 and confirm the
latest public law affecting chapter 9 is not newer than Public Law 119-93.
\item Read the full text of Public Law 119-93 and confirm it does not itself disturb
any of the ten conforming anchors.
\item Confirm 21 U.S.C. 360e-5 (FD\&C Act section 515D) is still unassigned in both the
Code and the FD\&C Act table of contents.
\item Re-read 21 U.S.C. 331 and confirm the last lettered prohibited act is still
(iii), so (jjj) remains open.
\item Confirm 21 U.S.C. 351 still runs through (j); section 355g through (b); section
360e(c)(1) through (H); section 360c through (k); section 360e-4 through (c); section
360k through (b).
\item Confirm the device-definition exclusion language in 21 U.S.C. 321(h) and the
clinical-decision-support exclusion at 21 U.S.C. 360j(o)(1)(A) through (E) are
unchanged.
\item Confirm every 510(k) cite reads 21 U.S.C. 360(k) and every preemption cite reads
21 U.S.C. 360k.
\item Confirm the conforming inserted-punctuation edits at 21 U.S.C. 360e(c)(1)(G) and
(H) still match the current text.
\item Record the verification date, the OLRC currency string, and the verifier of
record; attach the result to the frozen bill text before the sponsoring office files
via the eHopper.
\end{enumerate}

\uscprov{1}{\textbf{Sign-off.} The audit is complete only when a named human verifier
signs the result against the frozen introduction text and confirms, as of the filing
date, that the draft remains current through Public Law 119-93 and that section 515D,
codified at 21 U.S.C. 360e-5, is still unused. These lines are illustrative
placeholders.}

\sigline{Verifier of record, with OLRC currency string, date}

\sigline{Sponsoring office staff confirmation, date}
